\documentclass[a4paper,11pt]{article}

\usepackage[margin=2.5cm]{geometry}
\usepackage[T1]{fontenc}
\usepackage[utf8]{inputenc}
\usepackage{lmodern}
\usepackage{hyperref}

\title{Streaming Anomaly Detection Methods: Practical Considerations}
\author{sentinel-ebpf notes}
\date{\today}

\begin{document}

\maketitle

\section{Possible Models}
The first half of the methods are univariant, and cannot be use for our multivariant features.
Form the remaining half of them are supervised and cannot be used for our unsupervised setting.

remaining are:
\begin{itemize}
  \item k-Means [151]
  \item KNN [110]
  \item Torsk [60]
  \item EIF [58]
  \item iForest [83]
  \item HBOS [47]
  \item DBStream [55]
  \item CBLOF [59]
  \item COPOD [80]
  \item IF-LOF [36]
  \item LOF [22]
  \item COF [130]
  \item PCC [121]
\end{itemize}

\section{Context and constraints}
After this initial pre-filter (dropping univariate and supervised approaches), we further restrict ourselves to methods that:
\begin{itemize}
  \item operate on a \emph{very large, continuous event stream} (syscall-level events),
  \item tolerate an effectively huge ``window'' of normal behaviour,
  \item can be updated online with \emph{bounded state} and \emph{no constant full reconstruction},
  \item are reasonably simple to implement and maintain in production.
\end{itemize}

Below we group these remaining candidate methods into seven broad categories:


\section{Deep learning methods}
none of the remaining unsupervised multivariate candidates fall into this group.
\begin{itemize}
  \item Most are designed for batch or mini-batch training.
  \item Many focus on fixed-length time-series windows per entity (e.g. multivariate KPIs).
  \item Too complex for a baseline
\end{itemize}

\section{Local density / outlier detection methods}
Representative remaining methods: LOF, COF, IF-LOF, DBStream, CBLOF, COPOD, HBOS, iForest, EIF, KNN.

\subsection*{Neighbour-based local methods}
These methods fundamentally rely on \textbf{k-nearest neighbours} or small local neighbourhoods in the current window. For a large sliding window and high-dimensional features:
\begin{itemize}
  \item Exact kNN or LOF-style density requires either:
  \begin{itemize}
    \item full neighbour search over the window (too slow), or
    \item sophisticated \textbf{kNN and reverse-kNN indexes} that retain good performance in high dimensions.
  \end{itemize}
  \item Maintaining those indexes incrementally for hundreds of thousands of points per event type is a substantial engineering effort and still suffers from the curse of dimensionality.
  \item Even incremental variants (ILOF, DILOF, MCOD, DBStream) depend critically on such indexes and non-trivial incremental neighbourhood bookkeeping.
\end{itemize}

\subsection*{Closed-form / histogram-based variants}
Representative remaining methods: HBOS, COPOD, CBLOF/BLOF:
\begin{itemize}
  \item We already implement this flavour via our \textbf{Freq1D} baseline (per-feature histograms with decay and categorical handling).
\end{itemize}

\subsection*{Isolation Forest and variants}
Representative remaining methods:Isolation Forest (IF), Extended Isolation Forest (EIF):
\begin{itemize}
  \item Basic Isolation Forest variants are \emph{batch} methods in common libraries: they do not support true online updates and instead require periodic full or partial forest rebuilds to track drift.
  \item There are Incremental variants, but we already implement this flavour via our  Half-Space Trees (already implemented via River), 
  \item Subsequence-oriented variants (Subsequence IF/LOF) additionally target fixed-length time-series segments rather than per-syscall feature vectors.
\end{itemize}

\section{Data mining / subsequence methods}
none of the remaining unsupervised multivariate candidates fall into this group.

\section{Classic machine learning}
Representative remaining methods: k-Means, PCC.

\subsection*{Global projection and boundary methods}

k-Means and PCC are classic ML methods that optimise global objectives or low-rank structure. In a non-stationary stream:
\begin{itemize}
  \item Must adept to distribution drift, which typically means retraining or complicated incremental updates (or decaying the model).
  \item Online variants exist (e.g. incremental PCA, online SVM)
\end{itemize}

\subsection*{Clustering and tree ensembles}
Representative methods: Streaming k-Means
\begin{itemize}
  \item Cluster structure can become misaligned under complex non-stationary behaviour
  \item Assumes roughly spherical clusters
  \item Hard to tune over time.
\end{itemize}

\subsection*{Batch tree ensembles}
none of the remaining unsupervised multivariate candidates fall into this group.
\begin{itemize}
  \item primarily designed for batch training, no support for true streaming updates
  \item Expensive to rebuild
  \item Track drift
\end{itemize}


\section{Signal analysis methods}
Representative methods: Torsk.

\subsection*{Why they do not match our event model}
\begin{itemize}
  \item These methods assume a small number of well-structured continuous signals (e.g. KPI time-series), not sparse, heterogeneous event streams of high-dimensional feature vectors.
  \item They work on fixed-length windows per signal and require repeated transforms (FFT/DWT) over sliding windows, which is itself a form of constant state reconstruction for large windows.
  \item Applying them at syscall level would require defining and maintaining many per-entity time-series and transform states, adding significant complexity without clear benefit over our existing models.
\end{itemize}


\section{Stochastic sequence models}
none of the shortlisted methods are HMM/DBN-style sequence models

\end{document}